\section{Powers of bounded-degree bases have linear full-support MCA}
\label{app:power-family-mca}

A solution variety of quadratic degree need not produce quadratically
many exceptional labels on a received line. We prove this for powers of
bounded-degree polynomials, including a third-order differential equation
whose projective solution surface has genuinely quadratic degree.

Let $\mathcal D$ be a set of $n$ distinct field elements and let $f,g$ be
arbitrary words on $\mathcal D$. For a candidate $P$ of degree at most
$D$ and a label $z$, write
\[
 S(P,z)=\{x\in\mathcal D:f(x)+zg(x)=P(x)\}.
\]
Call $(P,z)$ bad at threshold $A>D$ if $|S(P,z)|\ge A$ and there are
no polynomials $F_0,G_0$ of degree at most $D$ agreeing with $f,g$,
respectively, on all of $S(P,z)$. Equivalently, $g$ has no degree-$D$
interpolant on this full support, since then $F_0=P-zG_0$ would work.
A label is bad if it has at least one bad candidate in the specified family.

\begin{theorem}[Fixed powers of bounded-degree bases]
\label{thm:power-family-linear-mca}
Fix $r\ge1$ and $0<\rho_-\le\rho_+<a\le1$. Consider candidates
\[
 P=cH^e,\qquad c\ne0,\quad H\text{ monic},\quad\deg H=r,
 \qquad D=re,
\]
where $\rho_-\le D/n\le\rho_+$ and $A\ge an$.
There are constants $C_{r,a}$, $c_{r,a}$, and
$n_0(r,a,\rho_-)$ such that, for $n\ge n_0$, every received line has
at most $C_{r,a}n$ bad labels if the characteristic is zero or
$p>c_{r,a}n$. In particular this holds along every sequence with
$p/n\to\infty$. The exponent $e$ is common to each candidate family.
\end{theorem}

The proof gives explicit, very large sufficient constants. It does not
cover arbitrarily varying root exponents or a base degree growing with
$n$, and it gives no improvement to the finite better.codes certificate.

\subsection{A high-power gcd estimate}

For a polynomial $H$, its divisor vector records the integer multiplicity
at each root over an algebraic closure. Affine independence of these
vectors is over $\mathbb Q$, not over the coefficient field.

\begin{lemma}\label{lem:power-family-gcd}
Let $t\in\{3,4\}$ and let $H_1,\ldots,H_t$ be monic degree-$r$
polynomials with affinely independent divisor vectors. Let $F,G$ be
independent homogeneous linear forms in $t$ variables, and put
\[
 N=\binom{m+t-1}{t-1},\qquad
 J=\binom{m+t-3}{t-3},\qquad M=N-J
\]
for an integer $m\ge1$. Suppose
\[
 e>\binom{N-1}{2}(tr-1),\qquad
 \operatorname{char}\F=0\quad\text{or}\quad p>rme.
\]
Outside the union $S$ of the roots of the $H_i$, the gcd of
$F(H_1^e,\ldots,H_t^e)$ and $G(H_1^e,\ldots,H_t^e)$ has degree at most
\begin{equation}\label{eq:power-family-K}
 K_t=\frac{(t-1)Jrme}{M}+\frac{(tr-1)(M-1)}2.
\end{equation}
The coefficient of $D=re$ is $4/(m+3)$ for $t=3$ and $18/(m+5)$
for $t=4$.
\end{lemma}

\begin{proof}
First, the $N$ degree-$m$ homogeneous monomials in the $H_i^e$ are
linearly independent. Otherwise take a minimal constant relation with
$q$ terms. Affine independence of the divisor vectors excludes $q=2$.
Any $q-1$ terms have a nonzero ordinary Wronskian: their degrees are
less than $p$, so a constant change of basis to distinct leading degrees
gives a nonzero Vandermonde leading coefficient. The same argument
applies in characteristic zero.

At each $s\in S$, omit a term of smallest order from this Wronskian.
Minimality makes all the resulting Wronskians constant multiples of one
another. If $A_0$ is the common gcd of the $q$ monomials, comparison of
root orders with the Wronskian degree gives
\[
 rme-\deg A_0\le\binom{q-1}{2}(|S|-1).
\]
But the left side is at least the rational degree of any ratio of two
distinct monomials. This ratio is an $e$-th power of a nonconstant
rational function, so its degree is at least $e$. Since $q\le N$ and
$|S|\le tr$, this contradicts the exponent assumption.

Let $V$ be the degree-$m$ part of the homogeneous ideal $(F,G)$.
Its dimension is $M$, because the quotient is a polynomial ring in
$t-2$ variables. Evaluation at $(H_1^e,\ldots,H_t^e)$ is injective by
the preceding paragraph. A basis $\psi_1,\ldots,\psi_M$ of the evaluated
space has a nonzero Wronskian $W$ with
\[
 \deg W\le Mrme-\binom M2.
\]
At $s\in S$, let $a_\alpha$ be the vanishing orders of the $N$
evaluated monomials. Intersecting a codimension-$J$ subspace with the
span of monomials of order at least each integer $k$ shows that it has
an adapted basis whose sum of orders is at least
$\sum_\alpha a_\alpha-J\max_\alpha a_\alpha$.
Indeed the dimension of each intersection loses at most $J$;
summing the resulting filtration dimensions proves this bound.
Cancellations can only increase the orders. A constant basis change
does not change the order of $W$, so
\[
 \operatorname{ord}_s W\ge
 \sum_\alpha a_\alpha-J\max_\alpha a_\alpha-\binom M2.
\]
Summing over $S$, use
$\sum_s\sum_\alpha a_\alpha=Nrme$ and
$\sum_s\max_\alpha a_\alpha\le trme$ to obtain
\[
 \sum_{s\in S}\operatorname{ord}_s W
 \ge Nrme-Jtrme-|S|\binom M2.
\]
The gcd $T$ of the two evaluated forms, with all factors supported on
$S$ removed, divides every $\psi_i$. The identity
$\operatorname{Wr}(Th_1,\ldots,Th_M)=T^M\operatorname{Wr}(h_1,\ldots,h_M)$
therefore contributes at least $M\deg T$ zeros outside $S$.
Combining the last two bounds and using $N-M=J$ proves
\eqref{eq:power-family-K}.
\end{proof}

This elementary argument is a specialization of the high-power gcd
phenomenon; for a broader characteristic-zero theorem see
Guo--Sun--Wang~\cite[Theorem 11]{guo-sun-wang-gcd}. The finite-characteristic
statement here follows directly from the displayed Wronskian argument.

\subsection{Dependent tuples and fixed-word lists}

Put
\[
 C_{t,r}=\sum_{j=1}^{t-1}\binom{(j+1)r-1}{r}.
\]
In a collection of $L$ labeled candidates with at most $T$ occurrences
of each monic direction $H$, at most $TC_{t,r}L^{t-1}$ ordered
$t$-tuples have affinely dependent divisor vectors. To prove this, take
the first dependent entry after $j$ preceding entries. Its roots lie in
the union of their at most $jr$ roots. There are at most
$\binom{(j+1)r-1}{r}$ monic degree-$r$ polynomials on this root set,
each with at most $T$ labels. Count the remaining entries freely and
sum over $j$. Repeated directions and repeated roots are included.

Consider a fixed word and candidates agreeing with it on at least
$bn$ positions, for a fixed $b>0$. For each direction $H$, deleting its
at most $r$ roots leaves disjoint agreement sets for distinct scalars.
Thus the scalar multiplicity is at most $n/(bn-r)$, and hence at most
$2/b$ for sufficiently large $n$.
An independent triple has at most $K_3+3r$ common agreement positions:
outside its root union apply Lemma~\ref{lem:power-family-gcd} to
$c_1Z_1-c_2Z_2$ and $c_1Z_1-c_3Z_3$.
Choose $m$ and then $n,e,p$ so that $K_3+3r\le b^3n/16$.
If $L\ge4/b$, and $b_x$ is the number of candidates agreeing at $x$,
convexity gives
\[
 \sum_x(b_x)_3\ge n(bL-2)^3\ge b^3nL^3/8.
\]
Here $(b_x)_3\ge\max(b_x-2,0)^3$ justifies the first inequality.
Independent triples contribute at most $b^3nL^3/16$, and dependent
triples at most $TC_{3,r}nL^2$. Consequently
\begin{equation}\label{eq:power-family-list}
 L\le L_0:=\max\{\lceil4/b\rceil,\lceil32C_{3,r}/b^4\rceil\}.
\end{equation}
This list estimate needs positive agreement, but not agreement above
the code rate.

\subsection{From fixed-word lists to full-support MCA}

Choose one bad candidate at each bad label and set $B=\lfloor A/2\rfloor$.
For a selected candidate $cH^e$, delete the roots of $H$ from its
support and group the remaining coordinates by $\lambda_x=g(x)/H(x)^e$.
Call the candidate heavy if some group has at least $B$ positions,
and light otherwise.

For a heavy group with ratio zero, $cH^e$ belongs to the fixed-word
list of $f$ at threshold $B$. Each fixed candidate gives at most $n$
bad labels: a bad support contains a coordinate with $g(x)\ne0$, and
that coordinate determines $z$.
For a heavy group with nonzero ratio, $\lambda H^e$ belongs to the
fixed-word list of $g$ at threshold $B$. Fix such a direction $H$ and
divide $f+zg$ by $H^e$ at nonroot coordinates. These are affine
functions of $z$. Except at at most $r$ labels at which a root coordinate
agrees accidentally, a bad support must contain two distinct such
functions. Otherwise their common two coefficients times $H^e$ give
the forbidden common witnesses on the full support. A support of at
least $A-r$ coordinates containing two different functions consumes
at least $A-r-1$ pairs of different functions. Each pair meets at at
most one label. This bounds the number of bad labels in this direction by
\[
 r+\left\lfloor\frac{\binom n2}{A-r-1}\right\rfloor.
\]
The list bound \eqref{eq:power-family-list}, with $b=a/3$ and
$B\ge A/3$, bounds both heavy contributions by $O_{r,a}(n)$.

In a light support, each ratio group has at most $B-1$ coordinates.
There are at least $(A-r)(A-r-B+1)/2$ pairs with unequal ratios.
For a fixed $H$, each such pair determines at most one pair $(z,c)$.
The number of selected light labels in one direction is therefore at most
\[
 T=\left\lfloor\frac{n(n-1)}{(A-r)(A-r-B+1)}\right\rfloor
 \le 8/a^2
\]
for sufficiently large $n$. Four candidates at distinct labels satisfy,
at each common agreement coordinate, the two independent forms
\begin{align*}
 &(z_2-z_3)c_1H_1^e+(z_3-z_1)c_2H_2^e+(z_1-z_2)c_3H_3^e=0,\\
 &(z_2-z_4)c_1H_1^e+(z_4-z_1)c_2H_2^e+(z_1-z_2)c_4H_4^e=0.
\end{align*}
For an independent quadruple their common support has size at most
$K_4+4r$, by Lemma~\ref{lem:power-family-gcd}.
Choose parameters so that this is at most $a^4n/32$.
If the total number $L$ of light labels is at least $6/a$, counting
ordered quadruples through coordinates gives
\[
 a^4nL^4/16\le (a^4n/32)L^4+TC_{4,r}nL^3.
\]
Thus the light labels number at most
\[
 L_1:=\max\{\lceil6/a\rceil,\lceil256C_{4,r}/a^6\rceil\}.
\]
Together the contributions are at most
\begin{equation}\label{eq:power-family-finite}
 L_0n+L_0\left(r+\left\lfloor\frac{\binom n2}{A-r-1}\right\rfloor\right)+L_1.
\end{equation}

For completeness, the parameters can be chosen in the required order.
Put $b=a/3$, choose $m_3\ge128/b^3-3$ and $m_4\ge1152/a^4-5$,
and let $\beta_t=(tr-1)(M_t-1)/2$ be the constant term of $K_t$.
Require
\begin{gather*}
 e>\max_{t=3,4}\binom{N_t-1}{2}(tr-1),\qquad
 p>D\max(m_3,m_4)\quad\text{in positive characteristic},\\
 \beta_3+3r\le b^3n/32,\qquad
 \beta_4+4r\le a^4n/64,\qquad n\ge6(r+1)/a,\quad A\ge3.
\end{gather*}
Since $D<n$ and $D\ge\rho_-n$, these hold for all sufficiently large
$n$ with a sufficiently large fixed lower bound on $p/n$. They ensure
every inequality used above. Equation~\eqref{eq:power-family-finite}
proves Theorem~\ref{thm:power-family-linear-mca}.
The zero candidate adds at most $n$ labels. The same coordinate-pair
argument handles all constant candidates, and a finite union over base
degrees $1,\ldots,r$ is also covered with the ambient degree cap unchanged.

\subsection{A quadratic-degree solution surface}

Fix $e\ge1$ in characteristic zero or $p>2e$. The nonzero polynomial
solutions of degree at most $2e$ to
\begin{equation}\label{eq:quadratic-power-ode}
 e^2P^2P'''+3e(1-e)PP'P''+(1-e)(1-2e)(P')^3=0
\end{equation}
are exactly $P=cH^e$ with $\deg H\le2$. Indeed, at a root of
multiplicity $m$, the coefficient of the lowest possible term, after
removing a nonzero local scalar, is $m(m-e)(m-2e)$. Thus every root
has multiplicity $e$ or $2e$. Conversely, substitution verifies every
$cH^e$; equivalently \eqref{eq:quadratic-power-ode} is the
denominator-cleared equation $(P^{1/e})'''=0$.
The monic base polynomial descends to the coefficient field by uniqueness.

The map $[H]\mapsto[H^e]$ from the projective plane of quadratics is
injective and base-point free. Its differential is injective on tangent
spaces since $e\ne0$ in the field. The image therefore has degree $e^2$.
One can also see the degree from the section
$P(u)=P(v)$, $P(w)=P(v)$ for three distinct points $u,v,w$.
There is no nonzero solution with $H(v)=0$, so normalize $H(v)=1$.
Then $H(u),H(w)$ independently range over the $e$-th roots of unity,
giving exactly $e^2$ distinct reduced points by quadratic interpolation.

This is a fixed third-order equation of total jet degree three with an
actual solution variety of degree $\Theta(n^2)$ at fixed positive rate.
Nevertheless Theorem~\ref{thm:power-family-linear-mca} bounds the bad
labels by $O(n)$ in its large-characteristic regime. The example therefore
does not establish a quadratic proximity-gap lower bound.

The supplementary exact checks in \path{research/power_family_mca/}
verify the classification on 156,194 monic polynomials, ten reduced
degree sections, seven Wronskian fixtures, 24 received lines, and
exhaustive small divisor-dependence counts. Twelve rational-arithmetic
fixtures check feasibility of the parameter choices without allocating
length-$n$ arrays. The finite checks supplement the proof; they do not
replace it or provide useful finite constants.
